\begin{textsample}{POS Dim 1 – vem\_ed\_17 – Score 13.00 – vem\_ed\_17/t080.txt}  \label{ex:f1_pos_132}
BOAS PRÁTICAS 2) Fatec Bragança Paulista e Uniminuto (Colômbia) – PCI “Lengua Española” envolve alunos de Gestão Empresarial e a \textbf{turma} de Licenciatura em Língua Espanhola conduzida por Margot Costas. “As \textbf{tarefas} \textbf{foram} pensadas a \textbf{partir} do perfil dos estudantes de ambas as instituições.
A primeira atividade \textbf{é} individual, a \textbf{tarefa} principal reúne os dois países e a última fica a cargo dos alunos colombianos, \textbf{que} ministrarão \textbf{pequenas} aulas de língua espanhola para os brasileiros.
Essa \textbf{é} uma \textbf{proposta} inovadora, \textbf{que} nunca fizemos em edições \textbf{anteriores}, por isso estamos otimistas com o resultado”, empolga-se Lilian.
Site do PCI: sites.google.com/view/eciaespanol/lengua-espa\%C3\%B1ola 3) Fatec Bragança e Universidad Distrital Francisco José de Caldas (Colômbia) – PCI “Entreculturas” \textbf{conta} com alunos brasileiros de Logística e colombianos \textbf{que} estudam Educação Infantil com o professor Miguel Nicholls.
Após as apresentações individuais em vídeo, há duas \textbf{tarefas} em equipes mistas, enfocando aspectos culturais de ambos os países. “Os alunos brasileiros praticam a língua espanhola de \textbf{forma} efetiva com os falantes nativos”, ressalta.
Site do PCI: sites.google.com/view/eciadistrital/entreculturas 4) Fatec Itu com DUOC (Chile) – PCI “Aspectos Sociales” une estudantes de Gestão Empresarial e de Avaliação de Projetos (professor Felipe Zambrano).
As \textbf{tarefas} ocorrem ao \textbf{longo} de um ano, algo inovador no planejamento de PCIs: no primeiro semestre, \textbf{serão} solicitadas atividades de rompehielo (quebra-gelo), \textbf{tarefa} intercultural e um levantamento das problemáticas sociais nas regiões em \textbf{que} os alunos vivem.
Isso prepara o cenário para a \textbf{tarefa} principal a \textbf{ser} elaborada no segundo semestre: um aplicativo sobre essas problemáticas.
Site do PCI: sites.google.com/view/eciachile/aspectos-sociales Sobre o \textbf{aprendizado} principal com esses projetos, Lilian sintetiza os aspectos mais \textbf{relevantes}: “O enfoque maior \textbf{é} a \textbf{interação} com os estrangeiros, de \textbf{forma} a propiciar aos alunos essa \textbf{vivência} além-fronteiras, usando a língua estudada em sala de aula.
Os aspectos culturais também \textbf{são} determinantes.
Percebeu-se, nesses anos de PCI, \textbf{quanto} os alunos se desenvolvem linguística e culturalmente com essa experiência”.

% matched lemmas: anterior, aprendizado, contar, forma, interação, longo, partir, pequeno, proposta, quanto, que, relevante, ser, tarefa, turma, vivência
\end{textsample}
